\section*{\centering Chapitre C3 : Changement d'état des fluides réels pures} 
\begin{enumerate}[label=\arabic{enumi} - , left=0pt, itemsep=1em] % Personnalisation de la numérotation
   
    \item Donner la définition de corps pur \par
    \begin{solution}
        ensemble d'entité chimiques identiques
    \end{solution}

    \item Donner la définition de phase \par
    \begin{solution}
        domaine dans lequel les variables intensives sont constantes
    \end{solution}

    \item Donner la définition de d'évolution ou d'équilibre d'un corps pur sous deux phases\par
    \begin{solution}
        si $\nu_1 \neq \nu_2$ : évolution vers une seule phase dans le sens des potenciels chimiques décroissant 
    \end{solution}

    \item Donner l'expression de la variation de l'entalphie massique\par
    \begin{solution}
        \[ \Delta h_{1 \rightarrow 2}(T) = h_2(T) - h_1(T)\]
    \end{solution}


    \item Donner la température et la variation de l'enthalphie massique de fusion et de vaporisation de l'eau à P=1bar\par
    \begin{solution}
        $T_{fus} = 0$ °C \\
        $T_{vap} = 100$ °C \\
        $\Delta h_{fus} = 334 kJ/kg$\\
        $\Delta h_{vap} = 2260 kJ/kg$\\

    \item Donner l'expression de la variation de l'entropie massique\par
    \begin{solution}
        \[ \Delta s_{1 \rightarrow 2} = \frac{\Delta h_{1 \rightarrow 2}}{T_{12}}\]
    \end{solution}
    \end{solution}


    \item Donner l'expression de s et h pour un mélange diphasé\par
    \begin{solution}
        \[ h = x_1h_1 + x_2h_2\]
        \[ s = x_1s_1 + x_2s_2\]
    \end{solution}

    \item Donner l'allure de la courbe,l'équation des droites dans les différents domaines (liquide,liquide + gazeux et gazeux) d'un diagramme entropique\par
    \begin{solution}
        Cf III.2
    \end{solution}

    \item Donner l'allure de la courbe,l'équation des droites dans les différents domaines (liquide,liquide + gazeux et gazeux) et différentes transformations d'un diagramme entropique\par
    \begin{solution}
        Cf III.2
    \end{solution}

    \item Donner l'allure de la courbe,l'équation des droites dans les différents domaines (liquide,liquide + gazeux et gazeux) et différentes transformations d'un diagramme des frigoristes\par
    \begin{solution}
        Cf III.3
    \end{solution}

    \item Donner l'allure de la courbe,l'équation des droites dans les différents domaines (liquide,liquide + gazeux et gazeux) et différentes transformations d'un diagramme de clapeyron\par
    \begin{solution}
        Cf III.4
    \end{solution}
\end{enumerate}